% !TITLE 沁园春·长沙
% !AUTHOR 毛泽东
% !DYNASTY 近代
% !GENRE 词
% !ORDER 4

\begin{poem}
独立寒秋，湘江北去，橘子洲头\textsuperscript{1}。
看万山红遍，层林尽染；漫江碧透，百舸\textsuperscript{2}争流。
鹰击长空，鱼翔浅底\textsuperscript{3}，万类霜天\textsuperscript{4}竞自由。
怅寥廓\textsuperscript{5}，问苍茫大地，谁主沉浮\textsuperscript{6}？

携来百侣\textsuperscript{7}曾游。忆往昔峥嵘岁月稠\textsuperscript{8}。
恰同学少年，风华正茂\textsuperscript{9}；书生意气，挥斥方遒\textsuperscript{10}。
指点江山\textsuperscript{11}，激扬文字\textsuperscript{12}，粪土当年万户侯\textsuperscript{13}。
曾记否，到中流击水\textsuperscript{14}，浪遏飞舟\textsuperscript{15}？
\end{poem}

\begin{annotations}
\notetext{1}{橘子洲：又名水陆洲，是湘江中一个狭长的小岛，在长沙城西。毛泽东青年时代常与友人在此游览。}
\notetext{2}{百舸：许多船只。舸（gě），大船。}
\notetext{3}{鱼翔浅底：鱼儿在清澈见底的水中游动，像鸟儿在空中飞翔。“翔”字把鱼游写得像在飞，突出水的清澈。}
\notetext{4}{万类霜天：千万种生物在霜降后的秋天里。“霜天”指深秋。}
\notetext{5}{寥廓：广阔高远。怅寥廓——面对广阔的天地，心生感慨。}
\notetext{6}{谁主沉浮：谁掌握着升降沉浮的命运？沉浮，指国家的兴衰、时代的走向。}
\notetext{7}{百侣：许多同伴、朋友。指毛泽东在湖南第一师范求学时期的同学和挚友。}
\notetext{8}{峥嵘岁月稠：不平凡的岁月多么丰富。峥嵘，山势高峻的样子，比喻不平凡；稠，多。}
\notetext{9}{风华正茂：风采才华正当旺盛之时。形容青年人的朝气。}
\notetext{10}{挥斥方遒：意气奔放，正强劲。挥斥，奔放；方，正；遒（qiú），强劲有力。}
\notetext{11}{指点江山：评论国家大事。江山指国家。“指点江山”意为纵论天下。}
\notetext{12}{激扬文字：用激昂的文字（文章、言论）来激浊扬清。}
\notetext{13}{万户侯：古代封赏食邑万户的侯爵，代指高官权贵。把这些达官贵人视为粪土一样不值钱。}
\notetext{14}{中流击水：在江心拍水（游泳）。也暗用祖逖“中流击楫”的典故——祖逖率军北伐，在江中敲击船桨发誓，表示收复中原的决心。}
\notetext{15}{浪遏飞舟：波浪几乎阻挡了飞驰的船只。遏（è），阻止。形容游泳时掀起的浪花之大。}
\end{annotations}
